\subsection{The mean value theorem}

2.2.1. Let \(x,y\) be in \(E\), and let \([x,y]\) be the closed segment joining these two points. Let moreover \(f\) be a map from a neighborhood of \([x,y]\) into the space \(F\), differentiable at every point of \([x,y]\). Then \(f(x)-f(y)\) belongs to the closed convex hull of the set of points \(Df(z).(x-y)\) for \(z\) in \([x,y]\).

2.2.2. Let U be a connected open subset of E, and let f be a map from U into
F, admitting a zero derivative at every point of U; then f is constant on U.

2.2.3. Let U be a convex open subset of E, and f a map from U into F, differentiable at every point of U. Given a continuous seminorm γ on F and a real number M ≥ 0, the following conditions are equivalent:

\begin{enumerate}
\def\labelenumi{(\roman{enumi})}
\item
  For every \(x\) in \(\mathbf{U}\) , one has \(\| Df(x)\|_{\gamma}\leqslant M\) .
\item
  For every \(x\) and every \(y\) in \(U\) , one has \(\| f(x) - f(y)\|_{\gamma} \leqslant M . \| x - y\| \) .
\end{enumerate}

2.2.4. Let U be a neighborhood of a point \(x_0\) of E and let \(f\) be a function defined in the complement of \(x_0\) in U, with values in F. Suppose that \(f\) has a derivative \(Df(x)\) at every point \(x\) of U, \(x \neq x_0\) and that the function \(x \mapsto Df(x)\) has a limit \(D_0\) as \(x\) tends to \(x_0\) . Then, if dim (E) \(\geqslant 2\) , \(f\) has a limit at \(x_0\) and the function \(f\) extended by continuity to all of U is differentiable with derivative \(D_0\) at \(x_0\) ; the same holds if dim (E) = 1 and one supposes that \(f\) has a limit at \(x_0\) .

